% ScholarWeft LaTeX export — "book" doc type.
% See document.tex for the full explanation of how this file is used
% (--include-in-header, available substitutions/hooks/markers).
%
% Unlike document/article, "book" keeps a dedicated title PAGE (matching
% book.docx/book.odt's own section-break-separated cover page) — see
% export_latex's is_book branch. Chapter-opener pages need NO extra code to
% suppress the header: LaTeX's `book` class already puts every \chapter's
% own first page on the `plain` page style (page number only), and
% DocumentCompiler.py already emits top-level headings as \chapter — this is
% standard, unmodified LaTeX book-class behaviour. Heading 1 = \chapter,
% Heading 2 = \section for this doc type (document/article have no chapters,
% so their Heading 1/2 are \section/\subsection instead).
%
% Per-chapter footnote AND figure numbering are both the `book` class's own
% NATIVE, unconditional behaviour — book.cls resets the footnote counter at
% every \chapter (plain "1, 2, 3", no chapter prefix) and gives figures a
% "chapter.N" \thefigure (e.g. "Figure 2.3") natively, with no custom code
% needed at all (confirmed empirically — a bare \documentclass{book} with no
% \counterwithin already does both). This is deliberately NOT tied to the
% "Restart footnote and figure numbering per chapter" export checkbox — for
% book, chapter-scoped numbering is how the class works, not an optional
% style choice, matching DOCX/ODT's own "Figure C.N" convention for exactly
% this doc type. (document.tex/article.tex DO tie figure numbering to that
% checkbox, since article-class has no equivalent native behaviour to fall
% back on — see SWTOKFIGURECOUNTER there. Footnote restart is never applied
% for document/article either way: article-class footnotes are naturally
% continuous, and that's the desired "non-book top-level sections don't
% restart" behaviour, so there's nothing to add.)
%
% Roman-numeral frontmatter mirrors roman_frontmatter — \frontmatter/
% \mainmatter are emitted as raw LaTeX blocks directly into the compiled
% markdown by DocumentCompiler.py, since they must sit in the BODY, not the
% preamble.

% SWEXTRACLASSOPTIONS:

% ── Customize this template ────────────────────────────────────────────
% Copy this file (e.g. into your own Export Templates folder) and edit any
% of the lines below to get your own fonts/margins/link color/spacing.
%
% Noto Serif here is a VARIABLE font: one file with a continuous `wght` axis
% (100–900) and no separate Bold file. LuaLaTeX exposes that axis through
% fontspec's Weight= key, so BoldFeatures={Weight=700} below is the font's
% REAL 700 instance — not a synthetic fake-bold. Weight=400, the axis
% default, is normal text. (fontspec under XeLaTeX can't address variable
% axes, so the \else branch keeps the old synthetic FakeBold, used only when
% a .tex is compiled by hand with xelatex.) Width=/Slant=/RawAxis= work the
% same way if you ever need them.
\ifluatex
  % ── Unicode GLYPH FALLBACK ──────────────────────────────────────────────
  % LuaLaTeX draws only glyphs the current font HAS. Noto Serif lacks
  % many symbols (→ U+2192, ⚙ U+2699, ✓ U+2713 …) and all EMOJI, which then
  % show as tofu boxes. Build a fallback chain from whichever fonts are
  % installed (each guarded by \IfFontExistsTF) and hand it to luaotfload: a
  % glyph the main font lacks is drawn from the first listed font that has
  % it. Ordered Serif-style symbols first (matched to Noto Serif),
  % then monochrome emoji, CJK, and broad catch-alls for other scripts.
  % NOTE: symbols fonts differ in vertical position — Apple Symbols draws →
  % near the BASELINE, STIX Two Math / Noto Sans Math centre it — so a math/
  % symbols font precedes Apple Symbols. Colour emoji fonts (Apple Color
  % Emoji, Segoe UI Emoji) CANNOT be used (no colour-bitmap support);
  % install monochrome "Noto Emoji" (listed below) or use twemojis.
  \def\swfallbacklist{}
  \IfFontExistsTF{STIX Two Math}{\xdef\swfallbacklist{\swfallbacklist "STIX Two Math:mode=node;",}}{}
  \IfFontExistsTF{Symbola}{\xdef\swfallbacklist{\swfallbacklist "Symbola:mode=node;",}}{}
  \IfFontExistsTF{Apple Symbols}{\xdef\swfallbacklist{\swfallbacklist "Apple Symbols:mode=node;",}}{}
  \IfFontExistsTF{Noto Sans Symbols 2}{\xdef\swfallbacklist{\swfallbacklist "Noto Sans Symbols 2:mode=node;",}}{}
  \IfFontExistsTF{Noto Emoji}{\xdef\swfallbacklist{\swfallbacklist "Noto Emoji:mode=node;",}}{}
  \IfFontExistsTF{Noto Serif CJK SC}{\xdef\swfallbacklist{\swfallbacklist "Noto Serif CJK SC:mode=node;",}}{}
  \IfFontExistsTF{Songti SC}{\xdef\swfallbacklist{\swfallbacklist "Songti SC:mode=node;",}}{}
  \IfFontExistsTF{Hiragino Sans GB}{\xdef\swfallbacklist{\swfallbacklist "Hiragino Sans GB:mode=node;",}}{}
  \IfFontExistsTF{WenQuanYi Zen Hei}{\xdef\swfallbacklist{\swfallbacklist "WenQuanYi Zen Hei:mode=node;",}}{}
  \IfFontExistsTF{Arial Unicode MS}{\xdef\swfallbacklist{\swfallbacklist "Arial Unicode MS:mode=node;",}}{}
  \IfFontExistsTF{DejaVu Serif}{\xdef\swfallbacklist{\swfallbacklist "DejaVu Serif:mode=node;",}}{}
  \directlua{luaotfload.add_fallback("swfallback", {\swfallbacklist})}
  % ────────────────────────────────────────────────────────────────────────
  \setmainfont{Noto Serif}[BoldFont={Noto Serif},BoldFeatures={Weight=700},
    RawFeature={fallback=swfallback}]
\else
  \setmainfont{Noto Serif}[BoldFont={Noto Serif},BoldFeatures={FakeBold=3.5}]
\fi
\usepackage[letterpaper,margin=1.3in]{geometry}  % page size + margins (try a4paper, or a per-side margin like "top=1in,bottom=1in,left=1.25in,right=1in")
\definecolor{swlinkcolor}{HTML}{467886}        % link/citation color, as a hex RRGGBB (no '#') — matches book.docx/.odt's Hyperlink style
% ── Obsidian callout boxes ──────────────────────────────────────────────────
% scripts/sw-callouts.lua wraps standard Obsidian callouts in this environment,
% passing the type's background colour (a 10% tint, as in Obsidian) and border
% colour. Needs tcolorbox; when it is absent the environment degrades to an
% unboxed block so the content still appears. Tune the padding/rule here.
\IfFileExists{tcolorbox.sty}{\usepackage{tcolorbox}\tcbuselibrary{breakable}}{}
\IfFileExists{changepage.sty}{\usepackage{changepage}}{}
\newif\ifswcalloutindent
\ifcsname adjustwidth\endcsname\swcalloutindenttrue\else\swcalloutindentfalse\fi
\ifcsname newtcolorbox\endcsname
  \newtcolorbox{swcalloutinner}[2]{%
    colback=#1, colframe=#2, boxrule=0.4pt, arc=1.5pt,
    left=6pt, right=6pt, top=4pt, bottom=4pt, boxsep=0pt,
    breakable, before skip=0.3cm, after skip=0.3cm,
  }%
  % Indent the whole box 1.5cm from each margin (matching the DOCX/ODT callout
  % styles' left/right indent), with 0.2cm above/below to set it apart.
  \newenvironment{swcalloutbox}[2]{%
    \ifswcalloutindent\begin{adjustwidth}{1.5cm}{1.5cm}\fi
    \begin{swcalloutinner}{#1}{#2}%
  }{%
    \end{swcalloutinner}%
    \ifswcalloutindent\end{adjustwidth}\fi
  }%
\else
  \newenvironment{swcalloutbox}[2]{}{}%
\fi
\usepackage{setspace}\setstretch{1.35}         % paragraph line spacing — a setspace MULTIPLE (1 = single); this is NOT the same number as the DOCX/ODT 1.15 line-spacing multiple (they measure from a different baseline reference), so tune it by eye.
\SetSinglespace{1.35}                          % FOOTNOTE/FLOAT line spacing — setspace forces footnotes and floats to SINGLE spacing regardless of \setstretch, so they look cramped next to the body; set this to match \setstretch (or lower). Affects footnotes & floats only, not body text.
\raggedbottom                                  % let pages end wherever content naturally does — see document.tex
                                                % for why (stretchy spacing otherwise fills to an even bottom
                                                % margin, at the cost of inconsistent-looking gaps). Switch to
                                                % \flushbottom for even bottom margins instead (the traditional
                                                % choice for a twoside book, at this same cost).
% ────────────────────────────────────────────────────────────────────────

% Title-page logo hook — empty by default. Redefine to add a journal/
% institution logo above the title, e.g.:
%   \renewcommand{\swTitleLogo}{\includegraphics[width=3cm]{/path/to/logo.png}\\[1em]}
\newcommand{\swTitleLogo}{}

% ── Arabic / non-Latin RTL scripts (wrapper-free) ───────────────────────
% See document.tex for the full rationale. In short: LuaLaTeX + babel
% `bidi=basic` gives correct Arabic letter-joining AND right-to-left
% ordering (XeLaTeX cannot — its macro-based alternatives crash with "TeX
% capacity exceeded" at the first footnote), and `onchar=ids fonts`
% auto-detects any Arabic-script run (including Arabic-derived "Ajami") with
% NO wrapper, so inline Arabic just works. The Arabic font is Scheherazade
% New at 120% (Arabic looks smaller than Latin at the same nominal size),
% with Noto Naskh Arabic as the fallback; Chinese is provided too (font
% selection only — CJK needs no bidi). The nested \IfFontExistsTF guards
% let a missing font degrade instead of breaking the export.
%
% Right indent for Arabic verse lines (babel maps LaTeX's "left" indents to
% the RTL right side) — the verse default is nearly flush. Raise it for a
% deeper indent.
\newlength{\swArabicVerseIndent}\setlength{\swArabicVerseIndent}{5em}
\newcommand{\swArabicVerseSetup}{\setlength{\leftmargini}{\swArabicVerseIndent}}

\ifluatex
  \makeatletter
  \@ifpackageloaded{babel}{}{\usepackage[bidi=basic,shorthands=off]{babel}}
  \makeatother
  \babelprovide[import,onchar=ids fonts]{arabic}
  \IfFontExistsTF{Scheherazade New}{%
    \babelfont[arabic]{rm}[Scale=1.2]{Scheherazade New}%
  }{\IfFontExistsTF{Noto Naskh Arabic}{\babelfont[arabic]{rm}[Scale=1.2]{Noto Naskh Arabic}}{}}%
  \babelprovide[import,onchar=ids fonts]{chinese}
  \IfFontExistsTF{Noto Serif CJK SC}{%
    \babelfont[chinese]{rm}{Noto Serif CJK SC}%
  }{\IfFontExistsTF{Songti SC}{\babelfont[chinese]{rm}{Songti SC}}{}}%
  \newcommand{\arabicfont}{\selectlanguage{arabic}\swArabicVerseSetup}
\else
  % XeLaTeX fallback (manual .tex compile only): no paragraph-level bidi, so
  % mixed inline runs come out reversed. LuaLaTeX is strongly preferred.
  \newfontfamily\swArabicFont[Script=Arabic]{Scheherazade New}
  \newcommand{\arabicfont}{\swArabicFont\swArabicVerseSetup}
\fi

% Bullet-list levels alternate bullet/dash instead of LaTeX's own default
% bullet/dash/asterisk/centered-dot cascade.
\usepackage{enumitem}
\setlist[itemize,1]{label=\textbullet}
\setlist[itemize,2]{label=\textendash}
\setlist[itemize,3]{label=\textbullet}
\setlist[itemize,4]{label=\textendash}

% Heading 1 (\chapter) / Heading 2 (\section) styles — edit these to change
% appearance (bold/size/spacing/etc.); deeper heading levels use LaTeX's own
% defaults. No number prefix by default.
%
% \titlespacing's before/after values are FIXED lengths (no "plus/minus"
% stretch component) deliberately — see document.tex for why: the class's
% own default heading spacing is stretchy, so a heading's gap from the text
% above it can end up much larger on one page than another with no visible
% cause, whenever the page-breaking algorithm stretches it to fill a page.
% \chapter's own "before" skip is moot (it always starts a fresh page) but
% titlesec still wants a value.
% Heading NUMBERING is decided by DocumentCompiler.py (number_headings +
% latexize_headings), which emits every heading as raw LaTeX: a numbered
% chapter becomes \chapter{...} (so LaTeX supplies the number and the chapter
% counter increments, which is what figure numbering "chapter.N" relies on),
% and every other heading becomes the starred form \section*/\subsection*/...
% carrying whatever literal number the compiler wrote (e.g. "1.1").
% DocumentCompiler.py also emits \setcounter{secnumdepth}{0} into the body so
% only chapters can print an automatic number — sections can never be
% double-numbered. There is therefore no titlesec/titletoc configuration here:
% LaTeX's own defaults style the headings.

\usepackage{fancyhdr}
% Running heads match book.docx/book.odt: LEFT (even) page = "Author Name –
% Short Title"; RIGHT (odd) page = the current chapter title. \leftmark holds
% the last \markboth's first argument; book.cls's own \chaptermark would make it
% the uppercased, numbered "CHAPTER N. TITLE", so redefine it to the bare title.
\renewcommand{\chaptermark}[1]{\markboth{#1}{}}
\pagestyle{fancy}
\fancyhf{}
\fancyhead[LE]{SWTOKAUTHOR{} -- SWTOKSHORTTITLE}
\fancyhead[RO]{\leftmark}
\fancyfoot[C]{\thepage}
\renewcommand{\headrulewidth}{0.4pt}
% Book class + twoside forces a blank left-hand page whenever \chapter (or
% \frontmatter/\mainmatter, both raw-block-inserted into the body) would
% otherwise start on an even page. Left un-styled that blank page still
% carries the running header and a page number, which reads as a rendering
% glitch since there's no content on it. Standard LaTeX fix: make the
% auto-inserted blank page use the empty page style.
\let\swOrigClearDoublePage\cleardoublepage
\renewcommand{\cleardoublepage}{%
  \clearpage
  {\pagestyle{empty}\swOrigClearDoublePage}%
}
% See document.tex for why \hypersetup{pdftitle=...} isn't set here, and for
% how colorlinks/linkcolor above still reaches pandoc's own hypersetup call.
